\section{全仕事の比較結果}
\begin{table*}[t]
\centering\small
\caption{全24条件．$T$は要約役を含む論理token，$M$は全モデル要求，$A$はAPI数．}
\label{tab:whole}
\begin{tabular}{lcrrrcrrr}
\hline
課題 & 基準成否 & 基準$T$ & $M$ & $A$ & 比較成否 & 比較$T$ & $M$ & $A$ \\
\hline
\multicolumn{9}{l}{群1：全保持対別呼出し要約} \\
写真 & 成功 & 135,628 & 16 & 99 & 成功 & 139,871 & 17 & 102 \\
精算 & 成功 & 211,110 & 24 & 25 & 成功 & 211,110 & 24 & 25 \\
返金 & 成功 & 102,087 & 13 & 13 & 成功 & 102,087 & 13 & 13 \\
CSV & 未達 & 358,638$+?_{1}$ & 21 & 53 & 成功 & 440,461 & 32 & 77 \\
\hline
\multicolumn{9}{l}{群2：別呼出し対任意欄の同時要約} \\
写真 & 成功 & 139,871 & 17 & 102 & 成功 & 136,201 & 16 & 99 \\
精算 & 成功 & 211,110 & 24 & 25 & 成功 & 211,110 & 24 & 25 \\
返金 & 成功 & 102,087 & 13 & 13 & 成功 & 102,087 & 13 & 13 \\
CSV & 成功 & 440,461 & 32 & 77 & 未達 & 327,694 & 21 & 80 \\
\hline
\multicolumn{9}{l}{群3：別呼出し対必須欄の同時要約} \\
写真 & 成功 & 139,871 & 17 & 102 & 成功 & 129,374 & 16 & 102 \\
精算 & 成功 & 211,110 & 24 & 25 & 成功 & 211,110 & 24 & 25 \\
返金 & 成功 & 102,087 & 13 & 13 & 成功 & 102,087 & 13 & 13 \\
CSV & 成功 & 440,461 & 32 & 77 & 未達 & 438,073 & 28 & 65 \\
\hline
\end{tabular}
\par\smallskip\footnotesize $+?_{1}$は使用量不明１要求．基準と比較は各群内で別々に実行した．
全保持と任意・必須欄の間の直接比較ではない．
\end{table*}

表\ref{tab:whole}に全24条件を示す．全保持対別呼出しの群では，成功は3/4対4/4であった．写真の別呼出しは次入力を16,425から8,708 tokenへ短縮したが，要約生成に18,455 tokenを使い，全仕事では135,628から139,871へ約3.13\%増えた．実行役だけなら14,212 token減っており，要約役を除くと費用差の符号が逆になる．両側成功の３課題合計も448,825対453,068で約0.95\%増えた．CSVは別呼出しだけ成功したが，全保持の終端１要求の使用量が不明で，全４課題の正確な費用差は未確定である．

任意欄の群は，別呼出し4/4対同時出力3/4成功だった．要約を求めた写真とCSVの２応答とも要約欄を出さず，履歴の要約置換は０件だった．同時出力は777,092 tokenで別呼出し893,529より約13.03\%少ないが，CSVの成功を失った差を含む．入力の圧縮や品質維持の効果と解釈しない．

必須欄の群では２機会とも要約を出し，形式と長さの検査後に採用された．写真は別呼出しと同じく成功し，全仕事は139,871から129,374 tokenへ10,497，約7.50\%減った．APIは両側102，モデル要求は17から16となった．実行役費用は同時出力側で7,958増えたが，別要約18,455の差と合わせて総量が減った．従って専用要約の費用全額が節約されたわけではない．

一方，CSVは同時出力で次入力が21,999から11,110 tokenへ短くなったものの，438,073 tokenで次要求の予算を予約できず未達だった．別呼出しは440,461で成功した．同時出力の作業時間は約1,429.695秒で，別呼出しの約1,304.294秒より長かった．要約欄の欠落や使用量不明による失敗ではないが，終端理由だけから要約の意味的欠落が原因とは確定できない．

必須欄の全４課題は880,644対893,529 token，両側成功３課題は442,571対453,068であった．成功損失１を含む全体を性能維持の削減とはしない．精算と返金は全群で要約が起動せず，方式間のtoken差も０だった．この２件を圧縮の成功例へ数えない．固定した採用基準は３群とも未達であった．

各群の最初の介入前の完全要求は写真14位置，CSV12位置で全て一致した．介入のない精算24位置と返金13位置も一致した．この照合は，後続の行動や費用の不変，個別の要因の因果効果までは示さない．全８枠とserverを閉鎖してから成否と費用を確認し，結果を見た条件変更や失敗分の補充実行は行っていない．
